% !TITLE 渡荆门送别
% !AUTHOR 李白
% !DYNASTY 唐
% !ORDER 5

\begin{poem}
渡远荆门\textsuperscript{1}外，来从楚国\textsuperscript{2}游。
山随平野尽，江入大荒\textsuperscript{3}流。
月下飞天镜\textsuperscript{4}，云生结海楼\textsuperscript{5}。
仍怜故乡水，万里送行舟。
\end{poem}

\begin{annotations}
\notetext{1}{荆门：荆门山，在今湖北宜昌附近，是古代蜀地和楚地的分界处。过了荆门，便离开四川进入了湖北平原。}
\notetext{2}{楚国：泛指今湖北、湖南一带，春秋战国时属楚国。}
\notetext{3}{大荒：广阔无边的原野。群山随着平原的出现而渐渐消失，长江汇入广阔的荒野中流淌。}
\notetext{4}{飞天镜：月亮倒映在水中，像一面从天飞下来的镜子。}
\notetext{5}{海楼：海市蜃楼。云气升腾，在空中结成了海市蜃楼般的幻景。"结"和"飞"两个动词让静景有了动感。}
\end{annotations}

\begin{appreciation}
这是李白出蜀之后的第一首名作，写于他乘船东下、经过荆门山时。那一年他二十五岁，刚刚告别蜀中的山水，第一次见到荆楚平原的辽阔。

渡远荆门外，来从楚国游——荆门山在今湖北宜昌附近，是蜀楚分界之处。李白从蜀中出来，过了荆门，便算是进入了楚地。一个从来的字，带出了多少少年人的意气风发：他不是为了谋生而漂泊，他是来游历的，来看这个世界的。

山随平野尽，江入大荒流——这是李白最经典的写景名句之一。船行出峡，两岸的群山渐渐退去，眼前是一望无际的平原；长江从峡谷中奔涌而出，汇入广袤的原野。随字和入字用得极好：山是随着船行而渐尽的，江是主动投入大荒的。一个被动，一个主动，写出了空间转换时的动态感。读这两句，你能感受到一种压迫感的解除——峡谷是封闭的、压抑的，而平原是开放的、自由的。李白从蜀中出来，不仅是地理上的离开，更是精神上的挣脱。

月下飞天镜，云生结海楼——月亮倒映在水中，像一面从天上飞下来的镜子；云气升腾，在海面上结成了楼阁。飞和结都是动词，让静态的景物有了生命。海楼即海市蜃楼，在平原地区的水面上，因光线折射而形成的幻影。李白把它们写进了诗里，不是当作奇景来炫耀，而是作为他眼中世界的一部分：这个世界本来就充满了神奇，只是需要一双发现的眼睛。

仍怜故乡水，万里送行舟——最后一句忽然收束，从壮阔的天地回到了柔软的心情。他怜爱的不是故乡的山，不是故乡的人，而是故乡的水——那从蜀中一路跟着他的江水。水是无情的，但诗人的眼中，这水有了情意，万里迢迢地护送着他的小舟。这种拟人不是刻意的修辞，而是人在离别时自然产生的感觉：当你离开一个地方，你会觉得那个地方的万物都在挽留你。
\end{appreciation}
